\chapter{Arithmetic compactness and metric variation}

\begin{lemma}[Bounded lattice bases]
\label{lem:bounded-lattice-bases}\notready
Fix $n\geq1$, $0<c\leq C$ and $\varepsilon>0$. Full lattices
$\Lambda\subset\mathbb R^n$ with covolume in $[c,C]$ and
$\|x\|\geq\varepsilon$ for $0\ne x\in\Lambda$ admit bases whose
vectors have norms bounded by a constant depending only on
$n,c,C,\varepsilon$.
\end{lemma}
\begin{proof}
Let $v_n$ be the volume of the unit ball and put
$R=3(C/v_n)^{1/n}$. A ball of radius $R/2$ has volume greater than
$\covol\Lambda$. It contains two distinct points differing by a lattice
vector: otherwise its translates into a half-open fundamental
parallelepiped would be disjoint almost everywhere, contradicting the
volume inequality. More explicitly, integrate over that parallelepiped
the integer-valued sum of indicators of all translates of the ball;
its integral exceeds the parallelepiped's volume, so its value exceeds
one at some point. Their difference is a nonzero lattice vector of norm
less than $R$.

Choose a shortest nonzero vector $u$, of length $\ell\in[\varepsilon,R]$.
It is primitive, for divisibility by an integer of absolute value greater
than one would give a shorter vector. Extend it to an integral basis.
Orthogonal projection to $u^\perp$ gives a full lattice $\Lambda'$
with covolume $\covol\Lambda/\ell\in[c/R,C/\varepsilon]$.
Every nonzero projected vector has length at least
$\sqrt3\ell/2$: lift it and subtract an integer multiple of $u$ so
the longitudinal component has absolute value at most $\ell/2$;
the resulting nonzero lattice vector has length at least $\ell$.
Induction in dimension therefore supplies a bounded basis of
$\Lambda'$. Lift each basis vector and again adjust its longitudinal
component into $[-\ell/2,\ell/2]$. Together with $u$ these lifts form
a basis of $\Lambda$, and their norms are bounded by the square root
of the squared projected bound plus $R^2/4$. In dimension one the
single generator has norm equal to covolume, at most $C$, starting
the induction.
\end{proof}

\begin{lemma}[Compact lattice frames]
\label{lem:mahler-frames}\uses{lem:bounded-lattice-bases}\notready
Under the hypotheses of the preceding lemma, there is a compact set
of nonsingular real $n\times n$ matrices containing a basis matrix of
every such lattice.
\end{lemma}
\begin{proof}
Let $M$ be the common bound for basis vectors. The set of matrices whose
columns have norm at most $M$ and whose absolute determinant lies in
$[c,C]$ is closed and bounded in $\mathbb R^{n^2}$, hence compact.
The positive lower determinant bound makes every matrix nonsingular,
and the preceding lemma puts a basis matrix of each lattice in this set.
\end{proof}

\begin{lemma}[Trace systole of normalized semistable bundles]
\label{lem:trace-semistable-systole}
\uses{def:hermitian-degree,lem:degree-rank-one,lem:integral-part-generic-subspace}
\source{bost-specht-manuscript:page-0011}\notready
If $A$ is semistable with $\degn(A)=0$, every nonzero $x\in A$
satisfies $\sum_{\sigma:K\hookrightarrow\mathbb C}\log\|x_\sigma\|\geq0$.
Its image in the actual trace lattice consequently has norm at least
$\sqrt{[K:\mathbb Q]}$.
\end{lemma}
\begin{proof}
Let $L=A\cap Kx$. Its generic rank is one, its quotient is torsion free,
and $x$ is a nonzero integral section of $L$. Semistability gives
$\degn(L)\leq0$. The line degree formula reads
\[
 D\degn(L)=\log\#(L/\Ocal x)-\sum_\sigma\log\|x_\sigma\|.
\]
The finite quotient has at least one element, so its logarithmic cardinality
is nonnegative. This proves the first assertion without assuming $L$ free.
For every positive $u$, $\log u\leq u-1$. Apply this to
$u=\|x_\sigma\|^2$ and sum over the $D$ embeddings. The left hand
sum is twice the nonnegative logarithmic sum. Hence
$D\leq\sum_\sigma\|x_\sigma\|^2=\|x\|_{\mathrm{tr}}^2$.
Both norms are nonnegative, giving the asserted systole. This is the
logarithmic form of the arithmetic--geometric mean step in the source.
\end{proof}

\begin{lemma}[Arithmetic compactness]
\label{lem:arithmetic-compactness}
\uses{def:hermitian-degree,lem:restriction-covolume,lem:mahler-frames,lem:trace-semistable-systole,lem:degree-maximum-continuity}
\source{bost-specht-manuscript:page-0011}\notready
For fixed $K$ and $a\geq1$, every sequence of normalized semistable
rank-$a$ Hermitian $\Ocal$-bundles has a subsequence which, after
$\Ocal$-module isomorphisms, has one fixed underlying module and
convergent metrics. The limit is normalized and semistable.
\end{lemma}
\begin{proof}
For $0\ne x\in A$, the free line $\Ocal x$ has degree
$-D^{-1}\sum_v w_v\log\|x_v\|_v$. Semistability and normalization
therefore give $\prod_v\|x_v\|_v^{w_v}\geq1$. The arithmetic--geometric
mean inequality, applied to $D$ positive numbers with each complex
norm repeated twice, gives $\sum_vw_v\|x_v\|_v^2\geq D$.
The underlying Euclidean lattice has systole at least $\sqrt D$ and
covolume $\Delta_K^{a/2}$. Compact lattice frames give integral bases
whose Gram matrices converge to a positive-definite matrix $G$.

For $\alpha\in\Ocal$, multiplication by $\alpha$ has operator norm at
most $\max_v|\alpha_v|$. Convergence of the positive-definite Gram
matrices bounds their condition numbers and hence bounds all entries
of these integral multiplication matrices. Choose a fixed integral
basis $\alpha_1,\ldots,\alpha_D$ of $\Ocal$. There are only finitely
many possibilities for the associated $D$ matrices, so pass to a
subsequence on which all are constant. These matrices fix the entire
$\Ocal$ action, and identify every module in the subsequence with one
fixed module. The limit is still a positive-definite metric.

This metric comes from the required archimedean metrics. Indeed, the
orthogonal idempotents of $K\otimes_\mathbb Q\mathbb R$ are real
linear combinations of the fixed multiplication matrices, and are
self-adjoint in every metric. These are closed matrix equations, so
they persist in the limit. The components remain orthogonal. On each
complex component multiplication by $i$ is orthogonal, another closed
equation, and thus the restriction is a Hermitian metric. Determinants
are continuous, so the covolume formula gives normalization. For each
fixed submodule its determinant norm varies continuously, and its
degree remains at most zero in the limit. Thus all semistability
inequalities persist.
\end{proof}

\begin{lemma}[Attainment of the maximum saturated degree]
\label{lem:maximum-saturated-degree}
\uses{lem:degree-finiteness}
\lean{BostConjecture.Arithmetic.HermitianMetric.exists_max_degree_submodule,BostConjecture.Arithmetic.HermitianMetric.exists_degree_eq_maxDegree,BostConjecture.Arithmetic.HermitianMetric.maxDegree_nonneg}
\leanok
\source{bost-specht-manuscript:page-0011,bost-specht-manuscript:page-0012}
For a finite projective $\Ocal$-module $A$ with a Hermitian metric $h$,
there is a saturated submodule $P\subseteq A$ such that
$\degn_h(Q)\leq\degn_h(P)$ for every saturated $Q\subseteq A$.
The maximum, denoted $\mathcal D(h)$, is nonnegative.
\end{lemma}
\begin{proof}
The zero submodule is saturated because $A$ is torsion free, and its
degree is zero. By degree finiteness the set of saturated submodules
of nonnegative degree is finite and nonempty. Choose a member of
largest degree in this finite set. Any saturated submodule outside
the set has negative degree and therefore has degree no greater than
the chosen member. This proves both assertions, including when $A$
has rank zero.
\end{proof}

\begin{lemma}[Uniform comparison of induced degrees]
\label{lem:metric-degree-comparison}
\uses{def:hermitian-degree,lem:determinant-contraction,lem:degree-metric-scaling,lem:maximum-saturated-degree}
\lean{BostConjecture.Arithmetic.HermitianMetric.abs_degree_submodule_sub_le_of_fiberNorm_le,BostConjecture.Arithmetic.HermitianMetric.abs_maxDegree_sub_le_of_fiberNorm_le}
\leanok
\source{bost-specht-manuscript:page-0012}
Let $h,g$ be Hermitian metrics on a finite projective $\Ocal$-module
$A$. Suppose $t\geq0$ and
\[
 e^{-t}\|x\|_{h,\sigma}\leq\|x\|_{g,\sigma}
 \leq e^t\|x\|_{h,\sigma}
\]
for every complex embedding $\sigma$ and every vector in its complex
fiber. For every integral submodule $P\subseteq A$ of generic rank $b$,
\[
 |\degn_g(P)-\degn_h(P)|\leq bt.
\]
Consequently, writing $a=\operatorname{rk}A$,
$|\mathcal D(g)-\mathcal D(h)|\leq at$.
\end{lemma}
\begin{proof}
The same norm comparisons hold on the induced metrics of $P$, since
its fiber injects into that of $A$. The identity from the metric
$e^th$ to $g$ is a contraction, where $e^th$ means that every norm is
multiplied by $e^t$. The determinant contraction inequality and the
rescaling formula give
$\degn_h(P)-bt\leq\degn_g(P)$.
The lower norm comparison says that the identity from $e^tg$ to $h$
is a contraction. It gives the reverse bound
$\degn_g(P)-bt\leq\degn_h(P)$, proving the absolute-value estimate.
Every submodule has $b\leq a$. Apply the estimate to a saturated
maximizer for $g$ and then to one for $h$; their existence follows
from the preceding lemma. The two resulting inequalities give the
claim for $\mathcal D$.
\end{proof}

\begin{lemma}[Local comparison of actual metrics]
\label{lem:metric-topology-comparison}
\uses{lem:metric-degree-comparison}
\lean{BostConjecture.Arithmetic.HermitianMetric.continuous_fiberNorm,BostConjecture.Arithmetic.HermitianMetric.eventually_all_fiberNorm_compare}
\leanok
\source{bost-specht-manuscript:page-0012}
Give the Hermitian metrics on a fixed finite projective module the topology
of pointwise convergence of their fiber forms. For every metric $h$ and
$t>0$, all sufficiently close metrics $g$ satisfy
$\|x\|_{g,\sigma}\leq e^t\|x\|_{h,\sigma}$ and
$\|x\|_{h,\sigma}\leq e^t\|x\|_{g,\sigma}$ for every $\sigma,x$.
\end{lemma}
\begin{proof}
Fix an embedding and a complex basis $e_1,\ldots,e_a$ of its fiber.
If $x=\sum_i x_i e_i$, then
$g(x,x)=\sum_{i,j}\overline{x_i}x_j g(e_i,e_j)$.
This finite sum proves joint continuity in $g,x$, with the fiber carrying
the topology from $h$; in particular the fiber norm is jointly continuous.
On the $h$-unit sphere the inequalities
$e^{-t}<\|x\|_g<e^t$ hold when $g=h$.
They define an open neighborhood of each pair $(h,x)$.
The unit sphere is compact, so finitely many product neighborhoods give
a single neighborhood of $h$ on which both inequalities hold throughout
the sphere. For $x\ne0$ apply these inequalities to
$x/\|x\|_h$ and use homogeneity; for $x=0$ both inequalities are equalities.
The fiber of dimension zero causes no exception, since it contains only
zero. Finally there are finitely many complex embeddings, so the intersection
of their neighborhoods is a neighborhood with the asserted simultaneous bounds.
This also shows that the form topology agrees with convergence of Gram
matrices in any complex bases.
\end{proof}

\begin{lemma}[Uniform degree gap for strict stability]
\label{lem:strict-stability-degree-gap}
\uses{lem:degree-finiteness,lem:metric-degree-comparison}
\lean{BostConjecture.Arithmetic.HermitianMetric.IsStable.exists_degree_gap,BostConjecture.Arithmetic.HermitianMetric.isStable_of_fiberNorm_compare}\leanok
\source{bost-specht-manuscript:page-0012}
If $h$ is normalized and strictly stable, there is $\delta>0$ such
that $\degn_h(P)\leq-\delta$ for every nonzero proper saturated
submodule $P$. In particular a normalized metric $g$ with uniform
norm distortion at most $t\geq0$ remains strictly stable whenever
$a t<\delta$, where $a=\operatorname{rk}A$.
\end{lemma}
\begin{proof}
Let $S$ be the set of nonzero proper saturated submodules of degree
at least $-1$. It is finite by degree finiteness. If $S$ is empty,
take $\delta=1$. Otherwise choose $P_0\in S$ of largest degree and
put $\delta=\min(1,-\degn_h(P_0))$. Strict stability and normalization
give $\degn_h(P_0)<0$, so $\delta>0$. For a member of $S$ the
maximality of $P_0$ gives the claimed bound; outside $S$ the degree
is less than $-1\leq-\delta$. The metric-comparison estimate then
gives $\degn_g(P)\leq-\delta+a t<0$. Since $g$ is normalized,
these are exactly the strict stability inequalities.
\end{proof}

\begin{lemma}[The maximum degree and openness of strict stability]
\label{lem:degree-maximum-continuity}
\uses{lem:degree-finiteness,lem:maximum-saturated-degree,lem:metric-degree-comparison,lem:metric-topology-comparison,lem:strict-stability-degree-gap}
\lean{BostConjecture.Arithmetic.HermitianMetric.exists_degree_eq_maxDegree,BostConjecture.Arithmetic.HermitianMetric.continuous_maxDegree,BostConjecture.Arithmetic.HermitianMetric.isOpen_isStable_normalized}\leanok
\source{bost-specht-manuscript:page-0011,bost-specht-manuscript:page-0012}
For every Hermitian bundle $A$, the number
\[
 \mathcal D(A)=\max\{\degn(B):B\subseteq A\text{ saturated}\},
 \qquad\degn(0)=0,
\]
exists. It is continuous in the metrics. The strictly stable normalized
metrics on a fixed module form an open subset of the normalized metrics.
\end{lemma}
\begin{proof}
Degree finiteness shows that the degrees at least zero come from a
finite set of submodules, including zero, so their maximum exists and
dominates every other degree. If two sets of norms differ by factors
between $e^{-t}$ and $e^t$, a rank-$b$ determinant norm changes by a
factor between $e^{-bt}$ and $e^{bt}$. Each normalized degree therefore
changes by at most $bt$, and $\mathcal D$ changes by at most $at$.
Convergence of positive-definite metrics gives such bounds with $t\to0$,
proving continuity, also along the fixed-module subsequences above.

At a normalized strictly stable metric every nonzero proper saturated
submodule has negative degree. Degree finiteness applied with lower
bound $-1$ gives a finite set of exceptional submodules; those outside
the set have degree less than $-1$. Choose $\delta>0$ such that every
exceptional nonzero proper submodule has degree at most $-2\delta$,
and $2\delta\leq1$. The same bound then holds for all nonzero proper
submodules. A metric change with $at<\delta$ keeps all their degrees
negative. Restricting to normalized metrics proves openness.
\end{proof}

